% !TEX program = xelatex
\documentclass[11pt]{article}
\usepackage{fontspec}
\usepackage{xeCJK}
\setmainfont{CMU Serif}
\IfFontExistsTF{Noto Serif CJK SC}{\setCJKmainfont{Noto Serif CJK SC}}{%
  \IfFontExistsTF{Source Han Serif SC}{\setCJKmainfont{Source Han Serif SC}}{%
    \setCJKmainfont{FandolSong-Regular}%
  }%
}%
\usepackage[margin=1in]{geometry}
\usepackage{amsmath}
\usepackage{amssymb}
\usepackage{booktabs}
\usepackage{siunitx}

\sisetup{table-number-alignment=center}

\setlength{\parindent}{0pt}
\setlength{\parskip}{0.35\baselineskip}
\setlength{\tabcolsep}{4pt}
\renewcommand{\arraystretch}{0.92}

\makeatletter
\renewcommand{\maketitle}{%
  \begin{center}
    {\LARGE \@title\par}
  \end{center}
  \vspace{-0.6\baselineskip}
}
\makeatother

\title{Idea报告草稿：MGR-SID}
\author{}
\date{}

\begin{document}

\maketitle

\section{研究背景}

生成式推荐将目标 item 的预测改写为离散 token 序列生成问题：先为每个 item 构造固定长度的语义 ID（Semantic ID, SID），再由大语言模型根据用户历史生成目标 item 的 SID。该路线能够自然复用语言模型的训练与解码范式，但也把一个新的问题推到前台：\textbf{item tokenizer 的结构会直接决定后续推荐模型是否容易学习、是否容易生成、以及是否容易在局部歧义区域做出稳定判断。}

近来的工作已经逐渐认识到，单纯依赖文本语义构造 SID 并不足够，collaborative information 应当进入 tokenizer。然而，已有 collaborative SID 方法大多仍然采用如下范式：
\begin{itemize}
  \item 将协同信息抽象成一个或几个 collaborative embedding；
  \item 进行全局融合、全局去噪或全局对齐；
  \item 再把这一统一表示送入整个 tokenizer。
\end{itemize}

这种处理方式在 flat representation 上是自然的，但对 hierarchical SID 未必足够精确。因为不同 SID level 本身就在承担不同角色：上层更像粗粒度 routing，中层更像 subgroup organization，底层更像 leaf-level discrimination。因此，一个更合适的问题可能不是“是否要把 collaboration 融进去”，而是：\textbf{不同 SID level 到底应该保留什么样的 collaborative structure。}

\section{当前仓库中的经验事实}

\subsection{全局 collision 存在，但不是主矛盾}

在当前已经检查过的 SID index 中，完整 SID collision 的确存在，但整体处于 sub-percent 水平。以当前 Industrial semantic baseline 为例，完整 SID collision 约为 $0.43\%$。这说明：
\begin{itemize}
  \item global collision 不是完全不存在；
  \item 但它不足以单独解释当前系统的主要性能瓶颈。
\end{itemize}

因此，当前问题不宜继续被简单叙述成“不同 item 经常量化到同一个 SID”，而更应被理解为：\textbf{hierarchical SID 在局部层级上的判别性仍然不足。}

\subsection{主要错误模式是 local leaf ambiguity}

我们对现有结果进行的结构诊断表明：
\begin{enumerate}
  \item 模型并非总是跑到完全错误的语义子树；
  \item 很多错误样本已经写对了第一层，甚至前两层 prefix；
  \item 最终错误主要发生在最后一层 leaf token 上。
\end{enumerate}

换句话说，当前系统的剩余困难更像是：
\begin{quote}
模型已经进入了正确或近似正确的 semantic prefix 子空间，但仍然无法在局部子树内部稳定地区分多个语义接近、文本相似、行为上却应分开的 item。
\end{quote}

因此，当前最值得研究的失败模式不是 global routing failure，而是 \textbf{local leaf ambiguity}。

\subsection{naive collaborative fusion 会导致结构性风险}

在此前的最小前端融合尝试中，我们已经看到：如果直接把浅层协同特征与文本表示做全局融合，再送入现有量化链路，会带来明显的结构风险，甚至导致 tokenizer collapse。这说明：
\begin{itemize}
  \item collaboration 的确不该被简单地视为一个“额外特征”；
  \item front-end collaborative injection 要想成立，必须回答：
  \begin{quote}
  collaborative structure 到底应该以什么形式进入 hierarchical SID 学习。
  \end{quote}
\end{itemize}

\section{Graph Pilot：为什么要从 graph 角度重新思考}

为了避免继续停留在“时间窗口长度不同”的浅层 probe，我们进一步把协同信息改写为 graph proxies，构造了三种 train-only graph view：
\begin{itemize}
  \item \texttt{coarse\_graph}：归一化无向协同图上的 1-hop 加权分；
  \item \texttt{mid\_graph}：归一化无向协同图上的 2-hop diffusion；
  \item \texttt{local\_trans}：最后一个历史 item 到 target 的有向转移图概率。
\end{itemize}

我们继续在现有错误桶上做 probe，比较：
\begin{itemize}
  \item ground-truth target 在某一 graph view 下的分数；
  \item 模型预测 SID 下候选 item 的分数；
  \item 统计 ``target-better rate''。
\end{itemize}

\begin{table}[t]
  \centering
  \caption{Graph pilot：不同 graph view 在不同错误桶上的 target-better rate 与 coverage。}
  \small
  \begin{tabular}{l l S[table-format=1.6] S[table-format=1.6] S[table-format=1.6] S[table-format=1.6]}
    \toprule
    数据集 & View & {All} & {Same-$l_1$} & {Same-$l_2$} & {Coverage} \\
    \midrule
    Industrial & \texttt{coarse\_graph} & 0.145204 & 0.273230 & 0.398148 & 0.397048 \\
    Industrial & \texttt{mid\_graph}    & 0.079029 & 0.144912 & 0.206790 & 0.920495 \\
    Industrial & \texttt{local\_trans}  & 0.057605 & 0.132743 & 0.209877 & 0.082599 \\
    \midrule
    Office & \texttt{coarse\_graph} & 0.154866 & 0.274123 & 0.416667 & 0.438076 \\
    Office & \texttt{mid\_graph}    & 0.099123 & 0.114035 & 0.154762 & 0.955720 \\
    Office & \texttt{local\_trans}  & 0.062486 & 0.177632 & 0.428571 & 0.097775 \\
    \bottomrule
  \end{tabular}
\end{table}

这个 pilot 给出了三个重要信号：
\begin{enumerate}
  \item \textbf{graph signals 也是 non-uniform 的。} \texttt{coarse\_graph} 是最稳的全局 prior，而 \texttt{local\_trans} 虽然 coverage 很低，但在 deepest ambiguity bucket 上可能很强。
  \item \textbf{naive 的 mid-scale graph 不够好。} 直接 2-hop diffusion 不是一个足够成熟的中尺度设计，这说明中尺度结构不能被简单理解为 global 与 local 的线性插值。
  \item \textbf{真正关键的也许不是窗口长度，而是 graph structure 与 graph frequency。}
\end{enumerate}

因此，下一阶段值得研究的问题变成：
\begin{quote}
\textbf{不同 SID level 应该保留什么 graph-structured collaborative information，而不同 graph view 又该如何分别去噪？}
\end{quote}

\section{Motivation}

结合当前仓库中的所有观察，我们认为更合适的叙事不再是：
\begin{itemize}
  \item ``给 tokenizer 加协同特征''
\end{itemize}
而应当转为：
\begin{quote}
\textbf{在 hierarchical SID 中，graph structure 应以 level-specific 的方式参与 tokenization learning。}
\end{quote}

更具体地说：
\begin{itemize}
  \item 上层 SID 可能更需要 broad、stable 的 collaborative organization；
  \item 中层 SID 可能更需要一种比 global 更局部、又比 transition 更稳定的中尺度结构；
  \item 底层 SID 可能更需要 local transition-sensitive distinction。
\end{itemize}

与此同时，不同 graph view 的噪声类型也并不相同：
\begin{itemize}
  \item coarse graph 需要 debias 和 smoothing；
  \item mid graph 需要更结构化的 operator，如 community extraction、spectral residual、band-pass filtering；
  \item local graph 需要 pruning、confidence filtering 和 recency-aware denoising。
\end{itemize}

因此，一个合理的 hierarchical SID tokenizer 不应把 collaboration 当作单一 embedding 全局复用，而应把它建模成：
\begin{quote}
\textbf{不同 level 需要保留不同图结构，并在量化训练中以不同方式受到这些图结构的监督。}
\end{quote}

\section{方法：MGR-SID}

\subsection{核心思想}

我们提出 \textbf{MGR-SID（Multiplex Graph-Regularized Hierarchical Semantic IDs）}。

其核心思想是：
\begin{enumerate}
  \item 构造一个 train-only 的 denoised multiplex graph bank；
  \item 为不同 SID level 学习不同 graph mixture；
  \item 在量化训练时，不仅保留 semantic structure，还保留各层被选中的 graph structure；
  \item 通过 graph-regularized quantization 学到更适合生成式推荐的 hierarchical SID。
\end{enumerate}

与 feature fusion 不同，MGR-SID 不仅仅把 graph 当成 item representation 的辅助来源，而是让 graph structure 直接参与 SID learning itself。

\subsection{Multiplex Graph Bank}

我们从 train-only 交互数据中构造三张 item graph：
\begin{itemize}
  \item \textbf{$G_{\text{coarse}}$}：去偏的无向协同图，用于描述 broad compatibility；
  \item \textbf{$G_{\text{mid}}$}：中尺度图，用于描述 neither-global-nor-local 的中层子结构；
  \item \textbf{$G_{\text{local}}$}：去噪后的有向 transition graph，用于描述短程 next-item preference。
\end{itemize}

这三张图并不是同一个图的不同参数版本，而是刻意对应不同 collaborative resolution 与不同 noise profile 的结构视图。

这里需要特别说明：当前方法只固定了三类 graph role，并\textbf{没有}提前锁死 $G_{\text{mid}}$ 的唯一实现。对第一版方法而言，$G_{\text{mid}}$ 只需要满足：
\begin{itemize}
  \item 比 $G_{\text{coarse}}$ 更局部；
  \item 比 $G_{\text{local}}$ 更稳定；
  \item 能为 same-prefix ambiguity 提供额外判别力。
\end{itemize}

因此，$G_{\text{mid}}$ 的当前候选可以是：
\begin{itemize}
  \item community-aware graph；
  \item band-pass graph；
  \item diffusion residual graph。
\end{itemize}

哪一个版本最终进入主方法，需要由后续 cheap pilot 决定，而不是在这份方向稿中被先验锁定。

\subsection{View-Specific Graph Purification}

不同图视图的失效模式不同，因此不能使用统一的去噪方案。我们对三张图分别做 purification：
\begin{itemize}
  \item 对 $G_{\text{coarse}}$：
  \begin{itemize}
    \item popularity debias；
    \item support thresholding；
    \item low-pass smoothing 或 low-rank cleanup。
  \end{itemize}
  \item 对 $G_{\text{mid}}$：
  \begin{itemize}
    \item community normalization；
    \item spectral residual extraction；
    \item diffusion residual cleanup。
  \end{itemize}
  \item 对 $G_{\text{local}}$：
  \begin{itemize}
    \item recency weighting；
    \item support pruning；
    \item confidence filtering。
  \end{itemize}
\end{itemize}

记净化后的图视图为：
\[
\tilde{G}_{\text{coarse}},\quad
\tilde{G}_{\text{mid}},\quad
\tilde{G}_{\text{local}}.
\]

这里列出的 purification 仍然是\textbf{设计空间}，不是已经冻结的唯一实现。当前已经确定的是原则：
\begin{itemize}
  \item coarse / mid / local 三类图必须分开去噪；
  \item view-specific denoising 是方法主张的一部分；
  \item 具体算子可以在实现阶段选择最小可行版本。
\end{itemize}

\subsection{Level-Wise Graph Allocation}

对于每个 SID level，我们引入一个可学习的 allocation module，在 semantic anchor 与三类 purified graph view 之间产生 level-specific weighting。

这里故意\textbf{不}提前锁定唯一的参数化形式。第一版实现完全可以使用简单的 gate 或 softmax，但这条 idea 的核心并不依赖某一个具体公式，而依赖如下结构判断：
\begin{quote}
hierarchical SID 的不同 level 解决的是不同判别问题，因此 graph supervision 的分配也应当按 level 分别学习。
\end{quote}

需要强调的是，MGR-SID 的主张不是：
\begin{quote}
Level 1 必须 coarse，Level 2 必须 mid，Level 3 必须 local。
\end{quote}

它真正的主张是：
\begin{quote}
\textbf{不同 SID level 解决的是不同结构判别问题，因此最优的 graph allocation 应当按 level 分别学习。}
\end{quote}

\subsection{Graph-Regularized Quantization}

这是 MGR-SID 与普通 graph feature fusion 最根本的区别。

在每一层量化训练中，我们不仅要求模型保 semantic consistency，还要求它保该层被分配到的 graph structure。也就是说，在某一层量化目标中：
\begin{itemize}
  \item 图上应该接近的 item，更容易在该层被分到一致或相近的离散结构；
  \item 那些不应过早 collapse 的 item，不会在浅层被粗暴压缩到一起。
\end{itemize}

概念上，每层目标都可以理解为：
\begin{itemize}
  \item 一个 semantic anchor term，用来保证语义稳定性；
  \item 一个 graph regularization term，用来约束该层量化结果对被选中 graph mixture 的结构保持能力。
\end{itemize}

但这里同样不应假装具体损失已经完全定稿。当前更准确的说法是：
\begin{quote}
\textbf{graph regularization 是训练信号的角色定义，而不是某一个已经冻结的 penalty 公式。}
\end{quote}

后续实现时，$\mathcal{L}_{\text{graph}}$ 可以具体化为邻接保持、局部结构保持、或跨层分离约束中的一种简洁版本；真正关键的是它必须让 graph structure 直接参与 SID learning objective，而不只是参与前端表示融合。

因此，graph 在这里不是 auxiliary feature，而是 \textbf{SID learning objective} 的一部分。

\subsection{Semantic Anti-Collapse Anchor}

由于我们已经知道 noisy collaborative injection 容易导致 tokenizer collapse，因此 MGR-SID 保留一个强 semantic anchor，包括：
\begin{itemize}
  \item semantic consistency 或 reconstruction；
  \item code usage health checks；
  \item 可选的 prefix-stability constraint。
\end{itemize}

这一部分的作用是：
\begin{itemize}
  \item 让 graph structure 修补 semantic SID 的盲点；
  \item 但不让 collaborative noise 直接破坏原有的 semantic hierarchy。
\end{itemize}

\section{为什么这条方法比简单 fusion 更强}

MGR-SID 的 stronger point 有三点。

\subsection{graph views 是显式的}

它不把 collaboration 压成一个统一 embedding，而是显式区分：
\[
G_{\text{coarse}},\quad G_{\text{mid}},\quad G_{\text{local}}.
\]

\subsection{denoising 是 view-specific 的}

它不再假设一次全局去噪就够，而是承认不同 graph views 的失败模式不同。

\subsection{graph structure 直接监督 quantization}

这点最关键。MGR-SID 的主贡献不是 ``better graph features''，而是：
\begin{quote}
\textbf{better hierarchical SID learning under level-specific graph constraints.}
\end{quote}

\section{与其他候选方向的区别}

\subsection{不同于 plain graph feature fusion}

后者仍然很容易被 reviewer 看成：
\begin{itemize}
  \item 多个 graph feature；
  \item 再加 level-wise gate；
  \item 本质仍是 feature engineering。
\end{itemize}

MGR-SID 则把 graph structure 推进到 tokenizer learning objective。

\subsection{不同于 full graph-defined partitioning}

另一种更激进的想法是：直接用 graph coarsening 或 graph partitioning 定义整棵 SID hierarchy。它的确更 radical，但也更难稳定，更难与现有 semantic SID 兼容。MGR-SID 选择的是更 feasible 的 middle path。

\subsection{不同于 post-hoc local repair}

像 ACLR-lite 那样的后处理方法可以验证 local repair headroom，但主干仍然停留在 inference-time refinement。MGR-SID 试图把 graph structure 更早地推进到 tokenizer learning itself。

\section{可证伪主张}

MGR-SID 不是一个不可证伪的大想法，而是有非常明确的验证标准：
\begin{enumerate}
  \item MGR-SID 应优于 semantic-only SID；
  \item MGR-SID 应优于 parameter-matched 的 graph feature fusion；
  \item MGR-SID 应优于 uniform graph regularization；
  \item learned level-wise graph allocation 应呈现 non-uniform pattern；
  \item 不同 graph views 应在不同 ambiguity bucket 上承担不同作用。
\end{enumerate}

如果这些现象不成立，那么这条路线应当降级为有趣的分析 insight，而不是完整 paper method。

\section{最小实验路线}

MGR-SID 的第一版实现不应过度膨胀，而应保持狭窄：
\begin{itemize}
  \item 一张 coarse graph；
  \item 一张选定的 mid graph；
  \item 一张 local graph；
  \item 一套最小的 level-wise graph allocation；
  \item 一套最小的 graph-regularized quantizer。
\end{itemize}

最关键的第一轮实验不应直接全量开跑，而应按如下顺序推进：
\begin{enumerate}
  \item 先在 graph pilot 中确定一个可信的 $G_{\text{mid}}$；
  \item 构建 graph bank 与 purification cache；
  \item 先实现并比较：
  \begin{itemize}
    \item semantic-only；
    \item graph feature fusion；
    \item uniform graph regularization；
    \item MGR-SID。
  \end{itemize}
  \item 只有在 Industrial 上得到正信号后，再做：
  \begin{itemize}
    \item swapped allocation；
    \item no-denoising control；
    \item Office generalization。
  \end{itemize}
\end{enumerate}

\section{总结}

当前仓库中的经验事实已经非常清楚：
\begin{itemize}
  \item global collision 不是主要问题；
  \item local leaf ambiguity 是更核心的失败模式；
  \item collaboration 的确应该进入 SID learning；
  \item 但 naive global fusion 会带来严重结构风险。
\end{itemize}

Graph pilot 进一步表明：
\begin{itemize}
  \item collaborative signal 的效用并不是单尺度的；
  \item 真正重要的也许不是时间窗口，而是 graph structure 与 graph frequency；
  \item 因此，hierarchical SID 更适合接受 graph-structured、level-specific、view-specific denoised supervision。
\end{itemize}

MGR-SID 正是在这一判断上提出的方法：它不再把 collaborative information 视为一个统一 embedding，而是把它显式建模为多张图，并将这些图以不同方式、不同层级地注入到 hierarchical quantization 中。如果后续实验能够证明它稳定优于 graph feature fusion 与 uniform graph regularization，并学出可解释的非均匀层级分配，那么这条路线将成为当前最有潜力的前端 tokenizer 主线之一。

\end{document}
